% SPDX-License-Identifier: CC-BY-4.0
\documentclass[10pt]{article}
\usepackage[margin=0.8in]{geometry}
\usepackage{longtable,booktabs,array}
\usepackage[T1]{fontenc}
\usepackage{lmodern}
\usepackage{textcomp}
\usepackage[hidelinks]{hyperref}
\setlength{\LTcapwidth}{\textwidth}
\renewcommand{\arraystretch}{1.08}
\newcolumntype{L}[1]{>{\raggedright\arraybackslash}p{#1}}
\title{Provenance of the Electrostatic Material Library}
\author{David Meeker}
\date{1 October 2026}
\begin{document}
\maketitle

\section*{The library}

\texttt{statlib\_catalog.dat} is an electrostatic material library for FEMM, 94 entries, built 1 October 2026. An electrostatic material is a relative permittivity, $\varepsilon_x$ and $\varepsilon_y$ (they differ only for the anisotropic 6H-SiC entry, whose c-axis lies along y), and a volume charge density, zero here. Electrostatics is the zero-frequency limit, so each dielectric takes the lowest-frequency value its source gives; the frequency column says which. The solid-insulation, insulating-oil and laminate values are shared with the current-flow library and live in its CSV files.

A convenience entry, ``Vacuum'', with $\varepsilon_r = 1$ by definition, sits above the folders, as ``Air'' does in the magnetics library; it needs no source and is not counted or tabulated below. The library is generated by \texttt{statlib\_src/\allowbreak build\_statlib.py}; \texttt{statlib\_src/\allowbreak catalog.csv} lists every other entry, and the table below is generated from it. \texttt{SOURCES.md} records each source document with its SHA-256 hash.

\section*{Liquids}

The electrostatic solver treats every material as a perfect insulator, and water is not one. Even ultrapure water at 25~\textdegree C, with $\sigma = 5.5\ \mu$S/m and $\varepsilon_r = 78.3$ (the current-flow library's entry, from ITU-R P.527-6 and Light et al.), relaxes free charge in $\varepsilon_0\varepsilon_r/\sigma = 0.13$~ms. In a field held longer than that, water behaves as a conductor and its permittivity does not enter. The water entries therefore sit in a folder named ``Water (high-frequency limit; DC: use current flow)'', and a problem with water in it belongs in the current-flow solver, which models conduction and permittivity together. A library entry carries no comment, so the folder name is the only warning a user sees.

The organic liquids with $\varepsilon_r$ above 10 (the alcohols, glycerol, acetone, formamide, nitrobenzene and 1,2-di\-chloro\-ethane) sit in ``Polar Organic Liquids (high-frequency limit)'' for the same reason. Polar solvents dissolve and dissociate ionic impurities, so they conduct too, by an amount that depends on purity; no held source gives a value. The low-permittivity organic liquids and the insulating oils are good insulators and can be used as they stand.

\section*{Check against the CRC Handbook}

The \emph{CRC Handbook of Chemistry and Physics} \cite{crcstat} covers 51 of the entries; the ``CRC'' column gives its value, table and the library's difference from it (for gases, the difference in $\varepsilon-1$, marked $^*$). 43 agree within 5\%, 5 within 10\%, and 3 differ by more. The gases agree within 1\% in $\varepsilon-1$, but CRC's values descend partly from the same NBS Circular 537, so they are not independent. The liquids are an independent check: CRC's are Wohlfarth's Landolt--B\"ornstein compilation (1991), evaluated at the library's temperature from CRC's fits where it tabulates another; water agrees within 0.3\%, and the other liquids within 2\% except 2-propanol ($-$5\%) and glycerol ($-$7\%). Its glass values are at 100~MHz, its inorganic-solid values often in the infrared, and its polymer values ``typical or average'', so differences of several per cent there are expected.

The larger differences: Si$_3$N$_4$ (Dorf's 12.7 against CRC's 4.2 for a film; both are far from the 7--8 usually quoted for silicon nitride, and neither is a reliable figure); Ta$_2$O$_5$ (Dorf's 28 against 24 for the $\beta$ phase); nylon 6,6 (Callister's 4.0 against 3.5 for polyamides generally; nylon absorbs water, which raises it). Silicon's CRC value is at 4.2~K and is not comparable.

\section*{Entries}

{\small
\begin{longtable}[l]{@{}L{1.6in}lL{0.7in}L{0.4in}L{1.25in}L{1.45in}@{}}
\toprule
\textbf{Entry} & \textbf{$\varepsilon_r$} & \textbf{Frequency} & \textbf{Src.} & \textbf{CRC} & \textbf{Note} \\
\midrule
\endhead
\bottomrule
\endlastfoot
\multicolumn{6}{@{}l}{\textit{Gases (20 C, 1 atm)}} \\
Helium & 1.000065 & static & \cite{nbs537} & 1.0000645 (6-221, +1\%$^*$) & (eps-1) x 10$^{6}$ = 65.0 \\
Hydrogen & 1.0002538 & static & \cite{nbs537} & 1.0002532 (6-221, +0\%$^*$) & (eps-1) x 10$^{6}$ = 253.8 \\
Oxygen & 1.0004947 & static & \cite{nbs537} & 1.0004941 (6-221, +0\%$^*$) & (eps-1) x 10$^{6}$ = 494.7 \\
Argon & 1.0005172 & static & \cite{nbs537} & 1.0005169 (6-221, +0\%$^*$) & (eps-1) x 10$^{6}$ = 517.2 \\
Air (dry, CO2-free) & 1.0005364 & static & \cite{nbs537} & 1.000536 (6-221, +0\%$^*$) & (eps-1) x 10$^{6}$ = 536.4 \\
Nitrogen & 1.000548 & static & \cite{nbs537} & 1.0005474 (6-221, +0\%$^*$) & (eps-1) x 10$^{6}$ = 548.0 \\
Carbon dioxide & 1.000922 & static & \cite{nbs537} & 1.0009217 (6-221, +0\%$^*$) & (eps-1) x 10$^{6}$ = 922 \\
\multicolumn{6}{@{}l}{\textit{Liquids / Water (high-frequency limit; DC: use current flow)}} \\
Water, 20 C & 80.37 & static & \cite{nbs514} & 80.1 (6-199, +0\%) & 20 C; Section A standard liquid \\
Water, 25 C & 78.54 & static & \cite{nbs514} & 78.36 (6-15, +0\%) & 25 C; Section A standard liquid \\
Water, 50 C & 70 & static & \cite{nbs514} & 69.88 (6-15, +0\%) & 50 C; from 78.54 at 25 C and the circular's coefficient a = 0.00200 per C, log10(eps) = log10(eps25) - a (t - 25); ITU-R P.527-6 eq. 8 gives 70.3 \\
\multicolumn{6}{@{}l}{\textit{Liquids / Polar Organic Liquids (high-frequency limit)}} \\
Methanol & 33.62 & static & \cite{nbs514} & 33 (6-199, +2\%) & 20 C; Section A standard liquid (32.63 at 25 C) \\
Ethanol & 24.3 & static & \cite{nbs514} & 24.85 (6-199, -2\%) & 25 C \\
1-Propanol & 20.1 & static & \cite{nbs514} & 20.52 (6-199, -2\%) & 25 C \\
2-Propanol & 18.3 & static & \cite{nbs514} & 19.26 (6-199, -5\%) & 25 C \\
Glycerol & 42.5 & static & \cite{nbs514} & 45.64 (6-199, -7\%) & 25 C \\
Acetone & 20.7 & static & \cite{nbs514} & 20.49 (6-199, +1\%) & 25 C \\
Formamide & 109 & static & \cite{nbs514} & 111 (6-199, -2\%) & 20 C \\
1,2-Dichloroethane & 10.65 & static & \cite{nbs514} & 10.42 (6-199, +2\%) & 20 C; Section A standard liquid (10.36 at 25 C) \\
Nitrobenzene & 35.74 & static & \cite{nbs514} & 35.6 (6-199, +0\%) & 20 C; Section A standard liquid (34.82 at 25 C) \\
\multicolumn{6}{@{}l}{\textit{Liquids / Low-Permittivity Organic Liquids}} \\
Diethyl ether & 4.335 & static & \cite{nbs514} & 4.2666 (6-199, +2\%) & 20 C \\
Benzene & 2.2811 & static & \cite{nbs514} & 2.2825 (6-199, -0\%) & 20 C; Section A standard liquid (2.274 at 25 C) \\
Toluene & 2.379 & static & \cite{nbs514} & 2.374 (6-199, +0\%) & 25 C \\
Hexane & 1.89 & static & \cite{nbs514} & 1.8865 (6-199, +0\%) & 20 C \\
Cyclohexane & 2.023 & static & \cite{nbs514} & 2.0243 (6-199, -0\%) & 20 C; Section A standard liquid (2.015 at 25 C) \\
Carbon tetrachloride & 2.238 & static & \cite{nbs514} & 2.2379 (6-199, +0\%) & 20 C; Section A standard liquid (2.228 at 25 C) \\
Chlorobenzene & 5.708 & static & \cite{nbs514} & 5.6895 (6-199, +0\%) & 20 C; Section A standard liquid (5.621 at 25 C) \\
\multicolumn{6}{@{}l}{\textit{Liquids / Insulating Oils}} \\
Capacitor oil & 2.2 & 60 Hz & \cite{dorf} & -- & 25 C \\
Pipe cable oil & 2.15 & 60 Hz & \cite{dorf} & -- & 25 C \\
Self-contained cable oil & 2.3 & 60 Hz & \cite{dorf} & -- & 25 C \\
Heavy cable oil & 2.23 & 60 Hz & \cite{dorf} & -- & 25 C \\
Transformer oil & 2.25 & 60 Hz & \cite{dorf} & -- & 25 C \\
Alkyl benzene & 2.1 & 60 Hz & \cite{dorf} & -- & 25 C \\
Polybutene pipe cable oil & 2.14 & 60 Hz & \cite{dorf} & -- & 25 C; dielectric constant given at 1 MHz \\
Polybutene capacitor oil & 2.22 & 60 Hz & \cite{dorf} & -- & 25 C; dielectric constant given at 1 MHz \\
Silicone fluid & 2.7 & 60 Hz & \cite{dorf} & -- & 25 C \\
Castor oil & 3.74 & 60 Hz & \cite{dorf} & -- & 25 C \\
C8F16O fluorocarbon & 1.86 & 60 Hz & \cite{dorf} & -- & 25 C \\
\multicolumn{6}{@{}l}{\textit{Solid Insulation / Ceramics, Glass and Stone}} \\
Alumina (Al2O3) & 8.5 & 60 Hz & \cite{dorf} & -- &  \\
Porcelain (mullite) & 8.2 & 60 Hz & \cite{dorf} & -- &  \\
Steatite & 5.5 & 60 Hz & \cite{dorf} & -- &  \\
Magnesium oxide (MgO) & 9.65 & 60 Hz & \cite{dorf} & 9.65 (12-48, +0\%) &  \\
Glass (soda lime) & 6.25 & 60 Hz & \cite{dorf} & 6.75 (12-58, -7\%) &  \\
Mica & 6.9 & 60 Hz & \cite{dorf} & 6.9 (12-48, +0\%) &  \\
Porcelain & 6 & 60 Hz & \cite{callister} & -- & 60 Hz; Dorf Table 55.2 gives 8.2 for mullite porcelain \\
Fused silica & 4 & 60 Hz & \cite{callister} & 4.1 (12-58, -2\%) & 60 Hz \\
Pyrex 7740 borosilicate glass & 4.6 & 1 MHz & \cite{pyrex} & 5 (12-58, -8\%) & 1 MHz, 20 C \\
Marble, dolomitic (USGS sample, New York) & 7.7 & 1 kHz & \cite{usgs} & -- & 1 kHz, dry; tan delta 0.0137; single sample \\
Marble, pink (USGS sample, Colorado) & 11.7 & 1 kHz & \cite{usgs} & -- & 1 kHz, dry; tan delta 0.0213; single sample \\
\multicolumn{6}{@{}l}{\textit{Solid Insulation / Polymers}} \\
Polyethylene (LDPE) & 2.3 & 60 Hz & \cite{dorf} & 2.3 (13-17, +0\%) &  \\
Polyethylene (MDPE) & 2.3 & 60 Hz & \cite{dorf} & 2.3 (13-17, +0\%) &  \\
Polyethylene (HDPE) & 2.35 & 60 Hz & \cite{dorf} & 2.3 (13-17, +2\%) &  \\
Cross-linked polyethylene (XLPE) & 2.3 & 60 Hz & \cite{dorf} & 2.3 (13-17, +0\%) &  \\
EPR & 3.25 & 60 Hz & \cite{dorf} & -- & dielectric constant 3.0-3.5 given at 1 kHz only (middle used) \\
Polypropylene & 2.25 & 60 Hz & \cite{dorf} & 2.3 (13-17, -2\%) &  \\
PTFE & 2 & 60 Hz & \cite{dorf} & 2.1 (13-17, -5\%) &  \\
Glass-reinforced polyester premix & 6.3 & 60 Hz & \cite{dorf} & -- &  \\
Thermoplastic polyester & 3.55 & 60 Hz & \cite{dorf} & -- & dielectric constant 3.3-3.8 given at 100 Hz (middle used) \\
Polycarbonate & 3.17 & 60 Hz & \cite{dorf} & 2.92 (13-17, +9\%) &  \\
Epoxy (mineral filler) & 5 & 60 Hz & \cite{dorf} & -- &  \\
Epoxy (silica filler) & 3.85 & 60 Hz & \cite{dorf} & -- &  \\
Silicone rubber & 3.65 & 60 Hz & \cite{dorf} & -- &  \\
Phenol-formaldehyde & 5.3 & 60 Hz & \cite{callister} & -- & 60 Hz \\
Nylon 6,6 & 4 & 60 Hz & \cite{callister} & 3.5 (13-17, +14\%) & 60 Hz \\
Polystyrene & 2.6 & 60 Hz & \cite{callister} & 2.6 (13-17, +0\%) & 60 Hz \\
Kapton HN polyimide film & 3.4 & 1 kHz & \cite{kapton} & -- & 1 kHz, 25 and 50 um film (3.5 for 75 and 125 um), ASTM D150 \\
Delrin 150 acetal (POM-H) & 3.7 & 60 Hz & \cite{delrin} & 3.8 (13-17, -3\%) & 60 Hz, 73 F, ASTM D150 \\
Mylar polyester film (PET) & 3.3 & 100 Hz & \cite{mylar} & 3.25 (13-17, +2\%) & 100 Hz, 25 C, read from Figure 6 (Mylar C, 25 um); about 3.28 at 1 kHz \\
Kapton FN 120FN616 (FEP-coated polyimide) & 3.1 & frequency not stated & \cite{kaptonfn} & -- & frequency not stated; ASTM D150 \\
Kapton FN 150FN019 (FEP-coated polyimide) & 2.7 & frequency not stated & \cite{kaptonfn} & -- & frequency not stated; ASTM D150 \\
Kapton FN 250FN029 (FEP-coated polyimide) & 3 & frequency not stated & \cite{kaptonfn} & -- & frequency not stated; ASTM D150 \\
PVC, rigid (Trovidur PN) & 3.2 & frequency not stated & \cite{trovidur} & 3.39 (13-17, -6\%) & frequency not stated \\
Kraft paper, oil-impregnated & 3.3 & frequency not stated & \cite{dorf} & -- & frequency not stated; approximately 3.1-3.5 (middle used); impregnant about 2.2, dried cellulose 6.5-10 \\
Natural rubber (polyisoprene) & 2.6 & 1 kHz & \cite{crcstat} & -- & 1 kHz, 27 C; CRC gives typical values (2.5 at 1 MHz) \\
Neoprene (polychloroprene) & 6.6 & 1 kHz & \cite{crcstat} & -- & 1 kHz, 25 C; CRC gives typical values (6.3 at 1 MHz) \\
\multicolumn{6}{@{}l}{\textit{Solid Insulation / Thin-Film Dielectrics}} \\
SiO2 film & 3.9 & 1 kHz & \cite{dorf} & -- &  \\
Si3N4 & 12.7 & 1 kHz & \cite{dorf} & 4.2 (12-48, +202\%) &  \\
Ta2O5 & 28 & 1 kHz & \cite{dorf} & 24 (12-48, +17\%) &  \\
HfO2 & 35 & 1 kHz & \cite{dorf} & -- &  \\
\multicolumn{6}{@{}l}{\textit{Circuit Board Laminates}} \\
FR-4 epoxy-glass & 4.35 & 1 MHz & \cite{kreith} & -- &  \\
Polyfunctional FR-4 & 4.35 & 1 MHz & \cite{kreith} & -- &  \\
Multifunctional FR-4 & 4.35 & 1 MHz & \cite{kreith} & -- &  \\
Polyimide-glass & 4 & 1 MHz & \cite{kreith} & -- &  \\
BT epoxy & 3.975 & 1 MHz & \cite{kreith} & -- &  \\
Cyanate ester-glass & 3.55 & 1 MHz & \cite{kreith} & -- &  \\
PPO epoxy & 4 & 1 MHz & \cite{kreith} & -- &  \\
PTFE-glass & 2.5 & 1 MHz & \cite{kreith} & -- &  \\
Alumina substrate & 8.9 & 1 MHz & \cite{kreith} & -- &  \\
\multicolumn{6}{@{}l}{\textit{Semiconductors}} \\
Silicon & 11.7 & static & \cite{ioffe} & 12.1 (12-48, -3\%) & 300 K \\
Germanium & 16.2 & static & \cite{ioffe} & 15.8 (12-48, +3\%) & 300 K \\
Gallium arsenide (GaAs) & 12.9 & static & \cite{ioffe} & 13.13 (12-48, -2\%) & static, 300 K \\
Indium phosphide (InP) & 12.5 & static & \cite{ioffe} & 12.61 (12-48, -1\%) & static, 300 K \\
Gallium nitride (GaN, wurtzite) & 8.9 & static & \cite{ioffe} & -- & static, 300 K \\
Silicon carbide, 6H (c-axis along y) & 9.66 / 10.03 & static & \cite{ioffe} & 9.66 (12-48, +0\%) & static, 300 K; 9.66 perpendicular to c, 9.66/0.9631 = 10.03 parallel to c \\
Diamond & 5.7 & static & \cite{ioffe} & 5.66 (12-48, +1\%) & 10$^{2}$-10$^{4}$ Hz \\
\end{longtable}
}

\begin{thebibliography}{99}
\bibitem{nbs537} A. A. Maryott and F. Buckley, \emph{Table of Dielectric Constants and Electric Dipole Moments of Substances in the Gaseous State}, NBS Circular 537, 1953, Sec.~2.2.
\bibitem{nbs514} A. A. Maryott and E. R. Smith, \emph{Table of Dielectric Constants of Pure Liquids}, NBS Circular 514, 1951.
\bibitem{dorf} R. C. Dorf (ed.), \emph{The Electrical Engineering Handbook}, 2nd ed., CRC Press, 2000, Ch.~55 (R. Bartnikas), Tables 55.1 and 55.2. Also the text of Ch.~55, p.~1370.
\bibitem{kreith} F. Kreith (ed.), \emph{The CRC Handbook of Mechanical Engineering}, CRC Press, 1999, Table 10A.1.7.
\bibitem{callister} W. D. Callister and D. G. Rethwisch, \emph{Materials Science and Engineering: An Introduction}, 10th ed., Wiley, 2018, Table 18.4, p.~683.
\bibitem{ioffe} Ioffe Institute, \emph{New Semiconductor Materials} database, ``basic parameters'' pages, \url{http://www.ioffe.ru/SVA/NSM/Semicond/}, accessed 1 October 2026.
\bibitem{kapton} DuPont, \emph{Kapton HN polyimide film}, technical data sheet.
\bibitem{kaptonfn} DuPont, \emph{Kapton FN polyimide film}, technical data sheet.
\bibitem{delrin} \emph{Delrin 150 NC010 acetal (POM-H)}, technical data (ASTM D150).
\bibitem{mylar} DuPont Teijin Films, \emph{Mylar Electrical Properties}, Fig.~6 (values read from the graph).
\bibitem{pyrex} Corning, \emph{Glass Silicon Constraint Substrates} (Pyrex 7740).
\bibitem{trovidur} R\"ochling, \emph{Trovidur PN grey} (PVC-U), technical data sheet.
\bibitem{usgs} G. R. Olhoeft, \emph{Tables of Room Temperature Electrical Properties for Selected Rocks and Minerals with Dielectric Permittivity Statistics}, USGS Open-File Report 79-993, 1979, Table 2.
\bibitem{crcstat} W. M. Haynes, ed., \emph{CRC Handbook of Chemistry and Physics}, 97th ed., CRC Press, 2016: Permittivity (Dielectric Constant) of Water at Various Frequencies (p.~6-15, from Archer and Wang, JPCRD 1983), of Liquids (6-199, C.~Wohlfarth, from Landolt-B\"ornstein IV/6, 1991), of Gases (6-221, from NBS Circular 537 and Harvey and Lemmon 2005) and of Inorganic Solids (12-48, from Young and Frederikse, JPCRD 1973); Dielectric Constants of Glasses (12-58); Dielectric Constant of Selected Polymers (13-17, ``typical or average values''). Source of the natural-rubber and neoprene entries; otherwise used as a check.
\end{thebibliography}

\noindent Copyright 2026 David Meeker. This document is licensed under the \href{https://creativecommons.org/licenses/by/4.0/}{Creative Commons Attribution 4.0 International License (CC BY 4.0)}.

\end{document}
